\BourbakiUnnumberedChapter{术语索引}{术语索引}

阿贝尔闭包：V，第77页。

阿贝尔扩张：V，第77页。

绝对半单自同态：VII，第42页。

不等式相加：VI，第2页。

伴随（通过……得到的扩张）：V，第11页。

代数闭包：V，第22页。

代数中的代数元素：V，第16页。

代数扩张：V，第17页。

代数数：V，第20页。

代数闭域：V，第20页。

代数相关族、子集：V，第106页。

代数不相交扩张：V，第112页。

代数自由扩张族：V，第115页。

代数自由族、子集：IV，第4页，V，第106页。

阿基米德有序群：VI，第35页，习题31。

阿廷定理：V，第65页。

阿廷-施赖埃尔定理：VI，第22页。

阿廷-施赖埃尔理论：V，第91页。

相伴元素：VI，第5页。

贝祖恒等式：VII，第2页。

布丹-傅里叶法则：VI，第42页，习题21。

域的特征指数：V，第7页。

环的特征：V，第2页。

特征子模：VII，第67页，习题16。

中国剩余定理：VI，第33页，习题25。

闭的（域，代数上）：V，第20页。

闭的（域，可分上）：V，第45页。

闭半直线：VI，第28页。

闭的（相对代数上）域：V，第19页。

（阿贝尔）闭包：V，第77页。

（代数）闭包：V，第22页。

（代数可分）闭包：V，第45页。

（完全）闭包：V，第5页。

（相对代数）闭包：V，第19页。

闭包（相对p-根）：V，第25页。闭包（相对可分代数）：V，第43页。形式幂级数的系数：IV，第24页。多项式的系数：IV，第1页。与群、幺半群结构相容（序关系）：VI，第1页。与环结构相容（序关系）：VI，第19页。与交换幺半群结构相容（预序关系）：VI，第3页。互补定向：VI，第29页。复合扩张：V，第12页。级数的复合：IV，第97页，习题15。Ω中的共轭类：V，第52页。共轭元素：V，第52页；共轭扩张：V，第52页。常数多项式、常数项：IV，第1页。内容：VII，第16页。互素元素：VI，第13页。TS(M)中的余积：IV，第50页。三次扩张：V，第10页。循环扩张：V，第85页。循环向量空间（关于u）：VII，第29页。分圆扩张：V，第81页。分圆多项式：V，第81页。

可分解模：VII，第23页。

一个多项式的（泛）分解代数：IV，第73页。

分解定理：VI，第11页。

戴德金定理：V，第27页。

扩张的（不可分）次数：V，第46页。

域上代数的次数：V，第10页。

代数元素的次数：V，第16页。

多项式的（全）次数：IV，第3页。

有理分式的（全）次数：IV，第20页。

扩张的（可分）次数：V，第31页。

高度为 \(\leqslant 1\) 的p-根扩张中的（偏）导子：V，第103页。

形式幂级数的（偏）导数：IV，第32页。

多项式的（偏）导数：IV，第6页。

笛卡尔符号法则：VI，第42页，习题21。

对角、可对角化的自同态集：VII，第40页。

可对角化、被对角化的代数：V，第28页。

直接基：VI，第29页。

方向向量：VI，第28页。

半直线的方向：VI，第28页。

首一多项式的判别式：IV，第81页。

多项式的判别式：IV，第83页。

元素序列的判别式：V，第47页。

（代数）不相交的扩张：V，第112页。

（线性）不相交的扩张：V，第14页。

TS(M)中的除幂：IV，第45页。

级数的除幂：IV，第96页，习题13。

整除关系：VI，第5页。

（欧几里得）除法：IV，第10页。

（最大公）因子：IV，第12页，VI，第8页。

元素的因子：VI，第5页。

模的（初等）因子：VII，第24页。

二重根：IV，第15页。

特征空间：VII，第30页。

特征值：VII，第30页。

模的初等因子：VII，第24页。

初等对称多项式：IV，第61页。

平展代数：V，第28页。

欧几里得引理：VI，第15页。

欧几里得整环：VII，第49页，习题7。

欧拉恒等式：VI，第8页，IV，第89页，习题6。

欧拉-拉格朗日定理：VI，第26页。

有理分式在原点的展开：IV，第30页，V，第39页。

形式幂级数的指数：IV，第39页。

（指数型）序列：IV，第91页，习题1。

域的扩张：V，第9页。

（代数恒等式的）扩张原理：IV，第18页。

有限扩张：V，第10页。

有限生成扩张：V，第11页。

n次型：IV，第2页。

形式幂级数：IV，第24页。

（广义）形式幂级数：IV，第38页。

（有理）分式域：IV，第19页。

弗罗贝尼乌斯同态：V，第4页，V，第92页。

伽罗瓦扩张：V，第57页。

扩张、多项式的伽罗瓦群：V，第58页。

伽罗瓦理论：V，第67页。

模的伽马代数：IV，第92页，习题2。

高斯整数：VII，第1页。

（生成的）扩张：V，第11页。

（生成的）伽罗瓦扩张：V，第57页。

（生成的）拟伽罗瓦扩张：V，第55页。

扩张的生成族：V，第11页。

元素的最大公因数（GCD）：VI，第8页。

多项式的最大公因数：IV，第12页。

主理想的最大公因数：VII，第3页。

半直线（闭的，开的）：VI，第28页。

半直线（相反的）：VI，第28页。

汉克尔行列式：IV，第90页，习题1。

 \(p\) -根元素：V，第24页。

 \(p\) -根式扩张的高度：V，第26页。

 \(\mathbf{K}(\mathbf{T})\) 中元素的高度：V，第148页，习题11。

埃尔米特插值公式：IV，第88页，习题13。

埃尔米特简化型：VII，第68页，习题21。

形式幂级数的齐次分量：IV，第25页。

不完全域：V，第7页。

不完全度：V，第170页，习题1。

不可分解模：VII，第23页。

未定元：IV，第1页。

线性映射的指标：V，第132页。

欧拉函数：V，第79页。

不可分元素：VII，第17页。

膨胀同态：V，第70页。

不可分次数：V，第46页。

整理想：VI，第6页。

线性映射的不变因子：VII，第22页。

模的不变因子：VII，第20页。

子模的不变因子：VII，第19页。

自同态的不变量（相似性）：VII，第32页。

不可约元素：VI，第17页。

不可约多项式：IV，第13页。

等压多项式：IV，第3页。

若尔当分解：VII，第43页。

若尔当分解（乘性的）：VII，第46页。

若尔当矩阵：VII，第35页。

库默尔扩张：V，第162页，习题7。

库默尔理论：V，第88页。

拉格朗日插值公式：IV，第16页。

运算（多项式的）：IV，第94页，习题9。

首项系数：IV，第3页。

元素的最小公倍数（LCM）：VI，第8页。

主理想的最小公倍数：VII，第3页。

莱夫谢茨原理：V，第117页。

有序群的字典序积：VI，第7页。

线性拓扑：IV，第28页。

线性不相交扩张：V，第13页。

线性拓扑化代数：IV，第28页。

形式幂级数的对数：IV，第40页。对数导数：IV，第41页。吕罗特定理：V，第149页，习题11。

麦克莱恩准则：V，第43页，V，第123页。极大有序域：VI，第25页。元素的极小多项式：V，第16页。自同态的极小多项式：VII，第29页。闵可夫斯基引理：VII，第50页，习题9。首一多项式：IV，第3页。单生成扩张：V，第11页。单项式：IV，第1页。重因子（无重因子的多项式）：IV，第14页。元素的倍数：VI，第5页。重根：IV，第15页。特征值的（几何）重数：VII，第30页。根的重数：IV，第15页。初等因子的重数：VII，第24页。

负元素：VI，第4页。

负的 \(n\) -向量：VI，第29页。

负基：VI，第29页。

牛顿多边形：V，第150页，习题2。

牛顿关系式：IV，第70页，IV，第75页。

幂零分量：VII，第45页。

正规基，正规基定理：V，第73页。

正规扩张：V，第53页。

范数：V，第47页。

开半直线：VI，第28页。

相反半直线：VI，第28页。

（广义）形式幂级数的阶：IV，第25页，IV，第38页。

根的阶：IV，第15页。

可序域：VI，第39页，习题8。

有序扩张：VI，第21页。

有序域：VI，第20页。

有序群、幺半群：VI，第1页。

有序环：VI，第19页。

向量空间的定向：VI，第29页。

定向仿射空间、向量空间：VI，第29页。

\(p\) -进整数（环）：V，第96页。

\(p\) -根闭包：V，第25页。

\(p\) -根元素：V，第24页。

\(p\) -根扩张：IV，第19页。

\(p\) -挠群：VII，第10页。

\(\pi\) -准素分量：VII，第8页。

\(\pi\) -准素模：VII，第7页。

环的完全闭包：V，第5页。

完全域：V，第7页。

特征p的完全环：V，第5页。

多项式，多项式代数：IV，第1页。

多项式函数：IV，第4页。

多项式律：IV，第94页，习题9。

多项式映射：IV，第57页。

多项式映射（齐次）：IV，第55页。

正基：VI，第29页。

正元素：VI，第4页。

正 \(n\) -向量：VI，第29页。

幂级数（广义形式）：IV，第24页，IV，第38页。

预序群、幺半群：VI，第3页。

准素扩张：V，第139页。

素域、子域：V，第1页。

本原元素（定理）：V，第40页。

本原单位根：V，第81页。

主分式理想：VI，第6页。

主理想整环：VII，第1页。

主理想环：VII，第49页，习题6。

有序群的（字典序）积：VI，第7页。

有序群的积：VI，第7页。

积定向：VI，第29页。

Puiseux定理：V，第150页，习题2。

纯基：V，第106页。

纯扩张：V，第106页。

纯子模：VII，第55页，习题7。

毕达哥拉斯域：VI，第39页，习题8。

二次扩张：V，第10页。

二次互反律（定理）：V，第166页，习题23。

拟伽罗瓦扩张：V，第53页。

欧几里得除法中的商：IV，第11页。

商定向：VI，第29页。

有理分式：IV，第19页。

有理函数：IV，第21页。

既约环：V，第34页。

正则代数：V，第140页。

正则扩张：V，第141页。

域在扩张中的相对代数闭包：V，第19页。

在扩张域中相对代数封闭的域：V，第19页。互素多项式：IV，第12页。欧几里得除法中的余数：IV，第11页。不可约元素的（完全）代表系：VII，第3页。限制同态：V，第60页，V，第70页。两个多项式的结式：IV，第76页。多项式的根：IV，第14页。单位根：V，第78页。符号法则：VI，第19页。

（绝对）半单分量：VII，第45页。（绝对）半单自同态：VII，第42页。半单自同态：VII，第41页。半单模：VII，第9页。可分代数：V，第119页。可分代数闭包（相对）：V，第44页。可分代数扩张：V，第36页。可分闭包：V，第44页。可分次数：V，第31页。可分元素：V，第39页。可分扩张：V，第121页。可分多项式：V，第37页。可分闭域：V，第45页。可分超越基：V，第136页。E上取值于F的级数：IV，第96页，习题11。（Γ-）集：V，第75页。元素的符号：VI，第20页。相似自同态、矩阵：VII，第29页。自同态的相似不变量：VII，第32页。单模：VII，第25页。单根：IV，第15页。分裂域、扩张：V，第21页。子扩张：V，第9页。子域：V，第1页。有理分式中可代入的元素：IV，第21页。形式幂级数中的代入：IV，第29页。多项式中的代入：IV，第4页。有理分式中的代入：IV，第21页。（初等）对称多项式：IV，第61页。对称形式幂级数：IV，第67页。对称多项式：IV，第61页。对称张量的对称积：IV，第43页。对称有理分式：IV，第67页。对称张量：IV，第42页。对称化张量：IV，第43页。

切线性映射：IV，第36页。形式幂级数的泰勒公式：IV，第31页。

多项式的泰勒公式：IV，第 8 页。形式幂级数的项、常数项：IV，第 24 页。多项式的项、常数项：IV，第 1 页。迹：V，第 47 页。超越基：V，第 109 页。超越次数：V，第 110 页。超越元：V，第 16 页。超越扩张：V，第 17 页。平移同态：V，第 70 页。可三角化自同态：VII，第 35 页。平凡扩张：V，第 9 页。

乌尔姆–齐平定理：VII，第 59 页，习题14。

幂单分量：VII，第 46 页。

幂单自同态：VII，第 46 页。

环的单位：VI，第 5 页。

首一多项式的通用分解：IV，第 73 页。

权：IV，第 3 页。

威尔逊公式：V，第 94 页。

多项式的零点：IV，第 14 页。
% semantic-fragments: legacy-input-seam
\relax{}
% Generated \tableofcontents replaces the OCR-captured printed contents.
